% GENERATED FILE. Edit canonical lessons, then run npm run generate:books.
% canonical-source-hash: fnv1a64:8b0d6f74ce7b135f
% canonical-lessons: TA-C476-kirumi, TA-C476-tatuppuci, TA-C476-mukakkavacam, TA-C476-cukataram, TA-C476-noyvayppatu, TA-R476-review-469-476

\chapter{Germs and Hygiene}
\label{ch:ta-germs-and-hygiene}
\hlchaptermodality{476}

\begin{chapteropening}
\textbf{By the end of this chapter:} \emph{I can talk about germs, vaccines, masks, hygiene and falling ill.}

Complete the last lesson of chapter 476: I can talk about germs, vaccines, masks, hygiene and falling ill.
\end{chapteropening}

\section[\texorpdfstring{\ta{கிருமி} (kirumi)}{கிருமி (kirumi)}]{\ta{கிருமி} (kirumi) — a germ; a microbe}

\label{lesson:TA-C476-kirumi}



\begin{quote}
Before the new one: say the Tamil for a scan, then the Tamil for a restricted diet.
\end{quote}

\subsection*{You'll want to know: \ta{கிருமி}}
\textbf{\ta{கிருமி}} — \emph{kirumi} — \textquotedblleft{}a germ; a microbe\textquotedblright{}.

\begin{quote}
\ta{கை} \ta{கழுவினால்} \ta{கிருமிகள்} \ta{போகும்}.
\end{quote}

\begin{quote}
\emph{kai kaḻuviṉāl kirumikaḷ pōkum.}
\end{quote}

\begin{quote}
If you wash your hands, the germs go away.
\end{quote}

Plural: \ta{கிருமிகள்}, \emph{kirumikaḷ}.

Its plural is regular: \ta{கிருமிகள்}.

\begin{cousinweb}
\textbf{\ta{கிருமி}} (\emph{kirumi}) is borrowed from Sanskrit \emph{kṛmi}, \textquotedblleft{}a worm, an insect\textquotedblright{} (Monier-Williams).
\end{cousinweb}

\begin{grammarlens}[title={What you've built}]
One more word for describing what you lived and felt.
\end{grammarlens}

\subsection*{Your turn}
\begin{itemize}
\raggedright
  \item \emph{Say it:} \emph{kirumi}
  \item \emph{Say it:} \emph{kirumi}, once more
  \item \emph{Say it:} \emph{kai kaḻuviṉāl kirumikaḷ pōkum.} — If you wash your hands, the germs go away.
  \item \emph{Recall:} say the Tamil for a scan, then the Tamil for a restricted diet, then say \emph{kirumi} again
\end{itemize}

\begin{tcolorbox}[breakable,colback=teal!4,colframe=teal!35!black,title={Before you move on}]
What does \ta{கிருமி} mean? (\textquotedblleft{}A germ; a microbe\textquotedblright{}.) Say it once more.
\end{tcolorbox}
\section[\texorpdfstring{\ta{தடுப்பூசி} (taṭuppūci)}{தடுப்பூசி (taṭuppūci)}]{\ta{தடுப்பூசி} (taṭuppūci) — a vaccine}

\label{lesson:TA-C476-tatuppuci}



\begin{quote}
Before the new one: say the Tamil for a restricted diet, then the Tamil for a germ.
\end{quote}

\subsection*{You'll want to know: \ta{தடுப்பூசி}}
\textbf{\ta{தடுப்பூசி}} — \emph{taṭuppūci} — \textquotedblleft{}a vaccine\textquotedblright{}.

\begin{quote}
\ta{குழந்தைக்கு} \ta{இன்று} \ta{தடுப்பூசி} \ta{போட்டார்கள்}.
\end{quote}

\begin{quote}
\emph{kuḻantaikku iṉṟu taṭuppūci pōṭṭārkaḷ.}
\end{quote}

\begin{quote}
The baby was vaccinated today.
\end{quote}

\ta{தடுப்பு} \textquotedblleft{}prevention\textquotedblright{} + \ta{ஊசி} \textquotedblleft{}injection\textquotedblright{}; \ta{தடுப்பூசி} \ta{போடு}, to vaccinate.

\begin{cousinweb}
\textbf{\ta{தடுப்பூசி}} (\emph{taṭuppūci}) is a Tamil compound: \ta{தடுப்பு}, from native \ta{தடு} \emph{taṭu}, \textquotedblleft{}to stop\textquotedblright{} (DEDR 3031), and \ta{ஊசி}, \textquotedblleft{}a needle\textquotedblright{}, borrowed from Sanskrit \emph{sūcī} (CDIAL 13551).
\end{cousinweb}

\begin{grammarlens}[title={What you've built}]
One more word for describing what you lived and felt.
\end{grammarlens}

\subsection*{Your turn}
\begin{itemize}
\raggedright
  \item \emph{Say it:} \emph{taṭuppūci}
  \item \emph{Say it:} \emph{taṭuppūci}, once more
  \item \emph{Say it:} \emph{kuḻantaikku iṉṟu taṭuppūci pōṭṭārkaḷ.} — The baby was vaccinated today.
  \item \emph{Recall:} say the Tamil for a restricted diet, then the Tamil for a germ, then say \emph{taṭuppūci} again
\end{itemize}

\begin{tcolorbox}[breakable,colback=teal!4,colframe=teal!35!black,title={Before you move on}]
What does \ta{தடுப்பூசி} mean? (\textquotedblleft{}A vaccine\textquotedblright{}.) Say it once more.
\end{tcolorbox}
\section[\texorpdfstring{\ta{முகக்கவசம்} (mukakkavacam)}{முகக்கவசம் (mukakkavacam)}]{\ta{முகக்கவசம்} (mukakkavacam) — a face mask}

\label{lesson:TA-C476-mukakkavacam}



\begin{quote}
Before the new one: say the Tamil for a germ, then the Tamil for a vaccine.
\end{quote}

\subsection*{You'll want to know: \ta{முகக்கவசம்}}
\textbf{\ta{முகக்கவசம்}} — \emph{mukakkavacam} — \textquotedblleft{}a face mask\textquotedblright{}.

\begin{quote}
\ta{மருத்துவமனையில்} \ta{எல்லோரும்} \ta{முகக்கவசம்} \ta{அணிந்திருந்தார்கள்}.
\end{quote}

\begin{quote}
\emph{maruttuvamaṉaiyil ellōrum mukakkavacam aṇintiruntārkaḷ.}
\end{quote}

\begin{quote}
Everyone in the hospital was wearing a face mask.
\end{quote}

\ta{முகம்} + \ta{கவசம்} \textquotedblleft{}armour\textquotedblright{}: the -\ta{ம்} drops and \ta{க} doubles.

\begin{cousinweb}
\textbf{\ta{முகக்கவசம்}} (\emph{mukakkavacam}) is a compound of native \ta{முகம்} \emph{mukam}, \textquotedblleft{}face\textquotedblright{} (DEDR 4889), and \ta{கவசம்}, \textquotedblleft{}armour\textquotedblright{}, borrowed from Sanskrit \emph{kavaca} (CDIAL 2957).
\end{cousinweb}

\begin{grammarlens}[title={What you've built}]
One more word for describing what you lived and felt.
\end{grammarlens}

\subsection*{Your turn}
\begin{itemize}
\raggedright
  \item \emph{Say it:} \emph{mukakkavacam}
  \item \emph{Say it:} \emph{mukakkavacam}, once more
  \item \emph{Say it:} \emph{maruttuvamaṉaiyil ellōrum mukakkavacam aṇintiruntārkaḷ.} — Everyone in the hospital was wearing a face mask.
  \item \emph{Recall:} say the Tamil for a germ, then the Tamil for a vaccine, then say \emph{mukakkavacam} again
\end{itemize}

\begin{tcolorbox}[breakable,colback=teal!4,colframe=teal!35!black,title={Before you move on}]
What does \ta{முகக்கவசம்} mean? (\textquotedblleft{}A face mask\textquotedblright{}.) Say it once more.
\end{tcolorbox}
\section[\texorpdfstring{\ta{சுகாதாரம்} (cukātāram)}{சுகாதாரம் (cukātāram)}]{\ta{சுகாதாரம்} (cukātāram) — hygiene; public health}

\label{lesson:TA-C476-cukataram}



\begin{quote}
Before the new one: say the Tamil for a vaccine, then the Tamil for a face mask.
\end{quote}

\subsection*{You'll want to know: \ta{சுகாதாரம்}}
\textbf{\ta{சுகாதாரம்}} — \emph{cukātāram} — \textquotedblleft{}hygiene; public health\textquotedblright{}.

\begin{quote}
\ta{பள்ளியில்} \ta{சுகாதாரம்} \ta{பற்றி} \ta{ஒரு} \ta{வகுப்பு} \ta{நடந்தது}.
\end{quote}

\begin{quote}
\emph{paḷḷiyil cukātāram paṟṟi oru vakuppu naṭantatu.}
\end{quote}

\begin{quote}
There was a class on hygiene at school.
\end{quote}

\ta{சுகாதாரமான}, hygienic.

\begin{cousinweb}
\textbf{\ta{சுகாதாரம்}} (\emph{cukātāram}) is usually traced to a learned compound of Sanskrit \emph{sukha}, \textquotedblleft{}ease, well-being\textquotedblright{}, and \emph{ādhāra}, \textquotedblleft{}support\textquotedblright{} (Monier-Williams).
\end{cousinweb}

\begin{grammarlens}[title={What you've built}]
One more word for describing what you lived and felt.
\end{grammarlens}

\subsection*{Your turn}
\begin{itemize}
\raggedright
  \item \emph{Say it:} \emph{cukātāram}
  \item \emph{Say it:} \emph{cukātāram}, once more
  \item \emph{Say it:} \emph{paḷḷiyil cukātāram paṟṟi oru vakuppu naṭantatu.} — There was a class on hygiene at school.
  \item \emph{Recall:} say the Tamil for a vaccine, then the Tamil for a face mask, then say \emph{cukātāram} again
\end{itemize}

\begin{tcolorbox}[breakable,colback=teal!4,colframe=teal!35!black,title={Before you move on}]
What does \ta{சுகாதாரம்} mean? (\textquotedblleft{}Hygiene; public health\textquotedblright{}.) Say it once more.
\end{tcolorbox}
\section[\texorpdfstring{\ta{நோய்வாய்ப்படு} (nōyvāyppaṭu)}{நோய்வாய்ப்படு (nōyvāyppaṭu)}]{\ta{நோய்வாய்ப்படு} (nōyvāyppaṭu) — to fall ill}

\label{lesson:TA-C476-noyvayppatu}



\begin{quote}
Before the new one: say the Tamil for a face mask, then the Tamil for hygiene.
\end{quote}

\subsection*{You'll want to know: \ta{நோய்வாய்ப்படு}}
\textbf{\ta{நோய்வாய்ப்படு}} — \emph{nōyvāyppaṭu} — \textquotedblleft{}to fall ill\textquotedblright{}.

\begin{quote}
\ta{மழையில்} \ta{நனைந்ததால்} \ta{அவன்} \ta{நோய்வாய்ப்பட்டான்}.
\end{quote}

\begin{quote}
\emph{maḻaiyil naṉaintatāl avaṉ nōyvāyppaṭṭāṉ.}
\end{quote}

\begin{quote}
He fell ill because he got soaked in the rain.
\end{quote}

Past (I …): \ta{நோய்வாய்ப்பட்டேன்}, \emph{nōyvāyppaṭṭēṉ}.

Written Tamil; in speech \ta{உடம்பு} \ta{சரியில்லை}.

\begin{cousinweb}
\textbf{\ta{நோய்வாய்ப்படு}} (\emph{nōyvāyppaṭu}) is a native fixed phrase built from \ta{நோய்} \emph{nōy}, \textquotedblleft{}disease\textquotedblright{} (DEDR 3793), \ta{வாய்} \textquotedblleft{}mouth\textquotedblright{} and \ta{படு} \textquotedblleft{}to undergo\textquotedblright{}.
\end{cousinweb}

\begin{grammarlens}[title={What you've built}]
One more word for describing what you lived and felt.
\end{grammarlens}

\subsection*{Your turn}
\begin{itemize}
\raggedright
  \item \emph{Say it:} \emph{nōyvāyppaṭu}
  \item \emph{Say it:} \emph{nōyvāyppaṭu}, once more
  \item \emph{Say it:} \emph{maḻaiyil naṉaintatāl avaṉ nōyvāyppaṭṭāṉ.} — He fell ill because he got soaked in the rain.
  \item \emph{Recall:} say the Tamil for a face mask, then the Tamil for hygiene, then say \emph{nōyvāyppaṭu} again
\end{itemize}

\begin{tcolorbox}[breakable,colback=teal!4,colframe=teal!35!black,title={Before you move on}]
What does \ta{நோய்வாய்ப்படு} mean? (\textquotedblleft{}To fall ill\textquotedblright{}.) Say it once more.
\end{tcolorbox}
\section[(dialogue)]{Looking back: from Enemies and Grudges to Germs and Hygiene}

\label{lesson:TA-R476-review-469-476}



\begin{quote}
Say the Tamil for enmity and hostility.
\end{quote}

\subsection*{Your turn}
Say these aloud:

\begin{itemize}
\raggedright
  \item enmity, hostility, a grudge, an enemy, to take revenge
  \item loyalty, betrayal, sacrifice, forgiveness, the cold shoulder
  \item a virtue, temperament, one’s nature, maturity, negligence
  \item a meek person, a naive person, a gullible person, a miser, a spendthrift
  \item charm, majesty, generosity, self-confidence, to admire
  \item a symptom, nausea, palpitations, to hiccup, to chafe
  \item treatment, to examine, the pulse, a scan, a restricted diet
  \item a germ, a vaccine, a face mask, hygiene, to fall ill
\end{itemize}

\begin{tcolorbox}[breakable,colback=teal!4,colframe=teal!35!black,title={Before you move on}]
Enmity? (\textbf{\ta{பகை}}, \emph{pakai}.) Hostility? (\textbf{\ta{விரோதம்}}, \emph{virōtam}.) A grudge? (\textbf{\ta{வஞ்சம்}}, \emph{vañcam}.) An enemy? (\textbf{\ta{எதிரி}}, \emph{etiri}.) To take revenge? (\textbf{\ta{பழிவாங்கு}}, \emph{paḻivāṅku}.) Loyalty? (\textbf{\ta{விசுவாசம்}}, \emph{vicuvācam}.) Betrayal? (\textbf{\ta{துரோகம்}}, \emph{turōkam}.) Sacrifice? (\textbf{\ta{தியாகம்}}, \emph{tiyākam}.) Forgiveness? (\textbf{\ta{மன்னிப்பு}}, \emph{maṉṉippu}.) The cold shoulder? (\textbf{\ta{பாராமுகம்}}, \emph{pārāmukam}.) A virtue? (\textbf{\ta{பண்பு}}, \emph{paṇpu}.) Temperament? (\textbf{\ta{சுபாவம்}}, \emph{cupāvam}.) One’s nature? (\textbf{\ta{இயல்பு}}, \emph{iyalpu}.) Maturity? (\textbf{\ta{பக்குவம்}}, \emph{pakkuvam}.) Negligence? (\textbf{\ta{அலட்சியம்}}, \emph{alaṭciyam}.) A meek person? (\textbf{\ta{சாது}}, \emph{cātu}.) A naive person? (\textbf{\ta{வெகுளி}}, \emph{vekuḷi}.) A gullible person? (\textbf{\ta{ஏமாளி}}, \emph{ēmāḷi}.) A miser? (\textbf{\ta{கஞ்சன்}}, \emph{kañcaṉ}.) A spendthrift? (\textbf{\ta{ஊதாரி}}, \emph{ūtāri}.) Charm? (\textbf{\ta{வசீகரம்}}, \emph{vacīkaram}.) Majesty? (\textbf{\ta{கம்பீரம்}}, \emph{kampīram}.) Generosity? (\textbf{\ta{தாராளம்}}, \emph{tārāḷam}.) Self-confidence? (\textbf{\ta{தன்னம்பிக்கை}}, \emph{taṉṉampikkai}.) To admire? (\textbf{\ta{மெச்சு}}, \emph{meccu}.) A symptom? (\textbf{\ta{அறிகுறி}}, \emph{aṟikuṟi}.) Nausea? (\textbf{\ta{குமட்டல்}}, \emph{kumaṭṭal}.) Palpitations? (\textbf{\ta{படபடப்பு}}, \emph{paṭapaṭappu}.) To hiccup? (\textbf{\ta{விக்கு}}, \emph{vikku}.) To chafe? (\textbf{\ta{உறுத்து}}, \emph{uṟuttu}.) Treatment? (\textbf{\ta{சிகிச்சை}}, \emph{cikiccai}.) To examine? (\textbf{\ta{பரிசோதி}}, \emph{paricōti}.) The pulse? (\textbf{\ta{நாடி}}, \emph{nāṭi}.) A scan? (\textbf{\ta{ஸ்கேன்}}, \emph{skēṉ}.) A restricted diet? (\textbf{\ta{பத்தியம்}}, \emph{pattiyam}.) A germ? (\textbf{\ta{கிருமி}}, \emph{kirumi}.) A vaccine? (\textbf{\ta{தடுப்பூசி}}, \emph{taṭuppūci}.) A face mask? (\textbf{\ta{முகக்கவசம்}}, \emph{mukakkavacam}.) Hygiene? (\textbf{\ta{சுகாதாரம்}}, \emph{cukātāram}.) To fall ill? (\textbf{\ta{நோய்வாய்ப்படு}}, \emph{nōyvāyppaṭu}.)
\end{tcolorbox}
